%! Polynomial Functions and Real Zeros — Unit 2, Lesson 2.3 — Student Packet Cover Sheet
\documentclass[10pt]{article}
\usepackage{precalculus-article}
\usepackage{precalculus-boxes}
\usepackage{ltablex}
\keepXColumns

\begin{document}

% Full-bleed plum banner
\begin{tikzpicture}[remember picture, overlay]
  \fill[plum] (current page.north west)
    rectangle ([yshift=-0.9in]current page.north east);
\end{tikzpicture}
\vspace{-0.8in}

\begin{center}
  {\color{white}\textbf{\LARGE Precalculus}}\\[4pt]
  {\color{lilacmid}\large\textbf{Unit 2: Polynomial Functions}}\\[6pt]
  {\color{white}\textbf{\large Lesson 2.3 \quad Polynomial Functions and Real Zeros}}
\end{center}

\vspace{0.22in}
\namedateperiod
\vspace{0.18in}

\begin{learningtargetbox}
\small
\begin{enumerate}[leftmargin=1.5em, itemsep=5pt]
  \item I can explain why \textit{zero}, \textit{root}, \textit{solution},
        and \textit{$x$-intercept} all point at the same numbers, and check
        whether a given value is a zero of a polynomial.
  \item I can use the factor/zero equivalence in both directions: read the
        real zeros off a factored polynomial, and write a factored polynomial
        from a list of zeros and a point on its graph.
  \item I can find the multiplicity of each real zero and predict whether the
        graph crosses the $x$-axis there or is tangent to it --- and match a
        pre-drawn graph to its factored form using that reasoning.
  \item I can find the degree of a polynomial from a table of equally spaced
        values by taking successive differences.
\end{enumerate}
\end{learningtargetbox}

\vspace{0.14in}

\begin{tocbox}
\small
\renewcommand{\arraystretch}{1.5}
\begin{tabularx}{\linewidth}{@{} c l X r @{}}
\rowcolor{lilac}
\textbf{\#} & \textbf{Component} & \textbf{Description} & \textbf{Score} \\
\midrule
1 & Warm-Up        & Spiral review of prerequisite skills                  & \blank{1.2cm} \\
2 & Guided Notes   & Zeros, factors, multiplicity, degree from differences & \blank{1.2cm} \\
3 & Group Activity & Four polynomials, two zeros: matching graphs to forms & \blank{1.2cm} \\
4 & Exit Ticket    & Independent check on multiplicity and degree          & \blank{1.2cm} \\
5 & Homework       & Practice and extension                                & \blank{1.2cm} \\
\midrule
  & \multicolumn{2}{r}{\textbf{Total}} & \blank{1.2cm} \\
\end{tabularx}
\end{tocbox}

\vspace{0.14in}

\begin{remindbox}
\small
Lesson 2.2 asked where a polynomial turns; today we ask where it \textbf{hits
zero}. One equivalence carries the whole lesson: $(x-a)$ is a factor exactly
when $a$ is a zero. Count how many times a factor repeats --- its
multiplicity --- and you can say, without graphing, whether the curve cuts
through the $x$-axis or just kisses it and turns back.
\end{remindbox}

\end{document}
